\section{Introduction}\label{sec:introduction}

Let $K$ be a field of characteristic zero. Keller's Jacobian conjecture asserts that a polynomial map $F\colon K^n\to K^n$ with $\det JF\in K^*$ has a polynomial inverse. Introduced by Keller~\cite{Keller}, the problem connects the geometry of affine space with the structure of polynomial algebras. For background on polynomial automorphisms, see~\cite{EssenBook}; for a recent account of the topology of polynomial maps, their nonproperness loci, and connections with the Jacobian conjecture, see El Hilany~\cite{ElHilany2025}.

The status of the general conjecture changed in July 2026 with Alp\"oge's three-dimensional counterexample. Tao~\cite{Tao2026} gives a geometric exposition of the example, while Gao~\cite{Gao2026} develops the tangent-sweep construction and its higher-dimensional extensions. The original map has constant Jacobian determinant $-2$ and identifies three distinct points. These developments show that the general Keller condition does not imply global invertibility in dimension at least three. They also require a distinction between that condition and the more restrictive nilpotent-Jacobian hypotheses studied here.

The reductions of Bass, Connell, and Wright~\cite{BCW} and Yagzhev~\cite{Yagzhev} reduce the general conjecture to maps $F=\id+H$ with $H$ cubic homogeneous and $JH$ nilpotent, with the dimension allowed to vary. These reductions can increase the dimension. Consequently, counterexamples to the general conjecture yield counterexamples within the reduced class in some dimension, but this argument does not determine the three-variable case with $JH$ nilpotent.

Recent work also examines the algebraic constraints imposed by the Jacobian condition. Lee and Li~\cite{LeeLi2024} use generalized Magnus expansions to obtain restrictions on the Newton polygons of inner polynomials associated with planar Jacobian pairs. Ram\'irez and Valqui~\cite{RamirezValqui2025} compute a Gr\"obner basis and describe the solution set of a polynomial system arising in a Laurent-series formulation of the planar problem. In a different direction, Hamada, Kato, and Komiya~\cite{HamadaKatoKomiya2025} relate a Tate-algebra version of the conjecture to the polynomial version under a Hausdorff adic-topology hypothesis. These approaches complement the study of normal forms, generic fibers, and explicit inverses pursued below.

Nilpotence also suggests a structural question. If $H(0)=0$ and $JH$ is nilpotent, must the components of $H$ be linearly dependent over $K$? The answer is affirmative in dimension two and, in any dimension, when the Jacobian has rank at most one; see~\cite{BCW,deBondtYan2014}. For homogeneous maps in dimension three, de Bondt and van den Essen proved the stronger conclusion of linear triangularizability~\cite{deBondtEssen}. The general dependence problem has counterexamples, including homogeneous counterexamples in higher dimensions~\cite{deBondt2006}. Thus a classification in dimension three must account for maps with linearly independent components.

Several classifications treat the case of independent components under additional hypotheses. Chamberland and van den Essen~\cite{ChamberlandEssen} considered maps
\[
 H=\bigl(u(x,y),v(x,y,z),h(u(x,y),v(x,y,z))\bigr).
\]
Yan and Tang~\cite{YanTang} studied restrictions on the degree in $z$, and Yan and de Bondt~\cite{YanBondt2019} treated related forms in arbitrary dimension. Classifications under further restrictions on variable dependence appear in~\cite{CastanedaEssen,Yan2022}. The manuscript of He and Yan~\cite{HeYan} studies maps $(u(x,y),v(x,y,z),h(x,y,z))$ under the inequality $\deg_zv\geq2\deg_zh$, and also treats the case $\deg_zh\leq3$. Its Problem~4.1 asks whether a common normal form remains valid without these restrictions.

This paper develops a classification of normalized three-variable maps with independent components and nilpotent Jacobian, and studies its consequences for tame and stably tame automorphisms. The argument is based on the family
\begin{equation}\label{eq:canonical-intro}
 \Ncal_P(x,y,z)=\bigl(P(y+x^2),\ z-2xP(y+x^2),\ P(y+x^2)^2\bigr),
 \qquad P\in K[t]\setminus K,\quad P(0)=0.
\end{equation}
Its Jacobian is nilpotent, and its components are linearly independent. Moreover, if
\[
 (X,Y,Z)=(x,y,z)+s\Ncal_P(x,y,z),
\]
then the identity
\begin{equation}\label{eq:intro-invariant}
 Y+X^2-sZ=y+x^2
\end{equation}
immediately gives a polynomial inverse. The main task is to reduce an arbitrary normalized map with independent components to this family.

The argument has two parts. First, Theorem~\ref{thm:dependent-pair} gives independent constant linear forms $\ell_1,\ell_2$ for which $\ell_1(H)$ and $\ell_2(H)$ are algebraically dependent. To find them, we pass from the image field to its relative algebraic closure in $K(x,y,z)$. Over this field, the generic curve is contained in an affine plane and satisfies an equation of degree at most two. Homogeneous triangularization restricts its directions at infinity. We exclude the line case and the conic with two distinct points at infinity; the remaining conic yields the dependent pair. A quadratic relation among the components then yields descent to $K$.

Second, Proposition~\ref{prop:pair-normal} converts a dependent pair into~\eqref{eq:canonical-intro}. After writing the pair as $f(r),g(r)$, we use a rational derivation with a slice to control the third component. A unit calculation rules out nonlinear dependence on a second invariant. Two commuting affine locally nilpotent derivations then show that $r$ is a polynomial in a coordinate $y+Q(x,z)$, where $\deg Q\leq2$. The remaining nilpotence identity forces the quadratic part of $Q$ to have rank one. Constant linear changes of coordinates complete the normalization. These changes are defined over $K$ and do not require the extraction of roots.

The resulting classification is stated in Theorem~\ref{thm:classification}. It leads to Theorem~\ref{thm1}: for $H\in K[x,y,z]^3$ with nilpotent Jacobian, the map $\id+H$ is tame. In the independent case we give an elementary factorization valid for the whole family $\id+sH$ over $K[s]$. In the dependent case, a constant linear conjugation makes one component vanish. A primitive relation over $K[z]$ then produces a polynomial change of basis in the remaining two coordinates. This gives a tame factorization over $K[s]$ as well.

We next consider maps in dimension $n\geq6$ of the form
\begin{equation}\label{eq:intro-block}
 \widetilde H=\bigl(H_1(x_1,x_2,x_3),\ldots,H_{n-3}(x_1,x_2,x_3),
 H_{n-2}(x),H_{n-1}(x),H_n(x)\bigr).
\end{equation}
Two three-variable nilpotent Jacobians occur as diagonal blocks of $J\widetilde H$. Theorem~\ref{thm2} gives their normal forms, using the fraction field of the base ring for the last block. This operation must be distinguished from conjugation of the full map by a matrix with variable coefficients. The constant term of the last block with respect to its three variables must also be subtracted.

The stable-tameness conclusion does not require descent of these rational changes of coordinates. We first show that the last block defines an automorphism over the polynomial base ring. Every residue-field specialization is tame by Theorem~\ref{thm1}. The residue-field criterion of Berson, van den Essen, and Wright~\cite[Theorem~4.12]{BEW} then gives stable tameness over that ring. Together with the tame base map and Suslin's theorem on polynomial matrices, this proves Theorem~\ref{thm3} for~\eqref{eq:intro-block}.

Section~\ref{sec:preliminaries} introduces the notation and records the algebraic results used below. Section~\ref{sec:dependent-pair} contains the dependent-pair argument. Section~\ref{sec:classification} gives the classification and tame factorizations. Section~\ref{sec:blocks} treats the higher-dimensional extensions. We conclude in Section~\ref{sec:conclusion} with several further questions.
